% INTRODUCTION 
\section{Introduction (1.5-2 pages)}


To infer statements about the dynamic response to exogeneous shocks, the estimation of the corresponding impulse response has relied on two fundamental approaches, either through vector autoregressive (VAR) or Local Projection (LP) models. Since the introduction of the latter, an extensive literature has emerged discussing the relative merits of the two. A first strand of literature has established that, under relatively general conditions, LPs and VARs estimate the same population impulse responses (Plagborg-Moller \& Wolf, 2021; Fusari et al. 2026). How-ever, in the discussion of finite sample behavior, LPs were often viewed as less sensitive to dy-namic misspecification whereas VARs were seen as potentially more statistically efficient (Li et al., 2024; Montiel Olea et al., 2025). 
\textcite{ludwig2026local} offers an analytical approach of explaining the observed finite sample bias-variance trade-off. He poses that LPs can be represented as a sequence of VAR estimations in-creasing in lag length as the estimated horizon increases and thus argues for a renewed thinking about how we can compare LPs and VARs – not by lag length, but by effective complexity. 

The first part of this paper aims to build on the discussion outlined above by asking whether conventional finite-sample conclusions are robust to alternative benchmark comparisons that relax the standard equal-lag restriction. To answer this question, the paper operationalizes Lud-wig’s (2026) central implication through a complexity-adjusted benchmarking exercise in which a fixed LP estimator is compared to progressively higher-order VAR specifications that proxy increasing effective dynamic flexibility (De Graeve \& Westermark, 2025). This comparison is conducted across empirically relevant sample sizes (see e.g. Herbst \& Johannsen 2021) to as-sess how the resulting bias–variance trade-off depends on the available data. 

The central research question is: Do LPs retain their apparent robustness advantage once they are compared not only to equal-lag VARs, but also to higher-order VARs that approximate the effective complexity implied by LP estimation? The contribution is not to establish a new LP–VAR equivalence result, but to ask whether standard Monte Carlo comparisons between LPs and equal-lag VARs are affected by the fact that LPs correspond to a more flexible dynamic benchmark.

The second part of the paper introduces controlled dynamic misspecification. While the baseline experiment studies complexity-adjusted LP–VAR comparisons under correctly specified linear dynamics, the main practical argument for LPs concerns their robustness when the estimated linear model is only an approximation to the true DGP (Montiel Olea et al., 2021). To study this channel explicitly, the paper introduces a controlled degree of misspecification into the DGP and examines how the relative performance of LPs, equal-lag VARs, and higher-order VARs chang-es as misspecification becomes more severe.

The simulation design proceeds in three steps.  First, we consider a correctly specified bivariate VAR(4) DGP, allowing LP(4), VAR(4), and higher-order VARs to be compared under common autoregressive information. Second, we introduce dynamic misspecification by adding a mov-ing-average shock component, thereby generating VARMA(4,1) dynamics while continuing to estimate finite-order VARs and LPs. Third, we introduce nonlinear misspecification through quadratic terms in the DGP while retaining linear estimators. This structure separates the role of dynamic approximation error from nonlinear functional-form misspecification and allows us to study whether conventional LP–VAR finite-sample rankings are robust to complexity-adjusted VAR benchmarks.

For each estimator, horizon, DGP, and sample size, we compute bias, variance, mean squared error, and empirical coverage rates. Bias and variance describe the finite-sample trade-off in point estimation, while coverage assesses whether the corresponding inference procedures re-main reliable under increasing dynamic or nonlinear misspecification.

